\documentclass[11pt]{article}

\usepackage[a4paper,margin=0.75in]{geometry}
\usepackage{amsmath,amssymb,amsfonts}
\usepackage{graphicx}
\usepackage{booktabs}
\usepackage{multirow}
\usepackage{array}
\usepackage{enumitem}
\usepackage{hyperref}
\usepackage{xcolor}
\usepackage{float}
\usepackage{caption}
\usepackage{subcaption}
\usepackage{setspace}

\setlength{\parskip}{0.5em}
\setlength{\parindent}{0pt}
\onehalfspacing

\hypersetup{
    colorlinks=true,
    linkcolor=blue,
    citecolor=blue,
    urlcolor=blue
}

\title{\textbf{MEL7118/ MCL775 Course Project Report}\\[0.5em]
\large Comparative Study of Model-Free Control Algorithms on Gridworld}
\author{
    Nirvan Nagrale \\
    2022ME12048 \\
    \and
    Kunal Parghi \\
    2022ME12043\\
    \and
    Aditya Yadav \\
    2022ME11320
}
\date{\today}

\begin{document}

\maketitle

\begin{abstract}
This report presents a comparative study of nine model-free reinforcement learning
(RL) algorithms implemented and evaluated on a $8\times8$ stochastic
Gridworld cost-minimisation task.  The algorithms span four families: off-policy
TD control (Q-learning, Double Q-learning, Entropy-Regularised Q-learning /
G-learning), on-policy TD control (SARSA, Expected SARSA), Monte-Carlo policy
gradient (REINFORCE, REINFORCE with baseline), and online actor-critic
(one-step AC, AC($\lambda$)).  All algorithms were trained for 1500 episodes
across 10 independent random seeds under identical environment instances,
discount factor ($\gamma = 0.85$), and a shared count-based learning-rate
schedule ($\alpha_t = n_t(s,a)^{-0.8}$).  Performance was measured via three
metrics: signed relative error (value-estimate bias), absolute relative error
(value-estimate accuracy), and policy evaluation error (quality of the induced
greedy policy relative to the known optimal $V^*$).  Results show that
Q-learning achieves the lowest final policy evaluation error ($\approx 0.045$
on a single representative seed), followed closely by G-learning and Double
Q-learning, while pure REINFORCE produces the largest error ($\approx 0.96$)
due to high-variance Monte-Carlo returns.  On-policy TD methods (SARSA,
Expected SARSA) converge more slowly than their off-policy counterparts in
this wall-heavy environment.  Entropy regularisation in G-learning provides
principled exploration annealing that improves upon standard $\varepsilon$-greedy
TD baselines.  Eligibility traces in AC($\lambda$) consistently improve over
one-step AC by propagating credit over multiple steps.  The findings confirm
classical RL theory: off-policy value-based methods with adaptive learning
rates outperform policy-gradient approaches on tabular, cost-minimisation MDPs
where Bellman signals are available at every step.
\end{abstract}

% ============================================================
\section{Introduction}
% ============================================================

Reinforcement learning (RL) is a sequential decision-making framework in which
an agent learns a behaviour policy by interacting with an environment, receiving
reward (or cost) signals, without being given a labelled dataset or an
explicit model of the world.  Its uniqueness lies in the \emph{credit-assignment
problem}: the agent must discover, purely from experience, which actions taken
at which states are responsible for eventual outcomes that may be separated from
the deciding moment by many time steps.

RL is the natural tool for a problem when three attributes are simultaneously
present.  First, the task is \emph{sequential}: decisions affect future states
and thus future decisions, so a greedy myopic policy is sub-optimal.  Second,
the \emph{environment model is unavailable or too complex to use}: transition
probabilities and cost functions are either unknown or have too many degrees
of freedom for explicit dynamic programming.  Third, there is a well-defined
\emph{scalar performance signal} that can be computed after each action.
Classic examples of use-cases in real lifeinclude robotic locomotion, game-playing, inventory management, and the navigation task studied in this project, where an
agent must reach a goal in a stochastic grid while minimising cumulative
movement cost and avoiding walls.

% ============================================================
\section{Problem Setup}
% ============================================================

\subsection{Gridworld Environment}

The environment is the $8 \times 8$ stochastic Gridworld introduced by Fox
et~al.~(2016).  States are cells of the grid, indexed by a flat integer
$s = 8r + c$ where $r \in \{0,\ldots,7\}$ is the row and $c \in \{0,\ldots,7\}$
is the column.

\paragraph{State space.}
Of the 64 raw grid cells, 15 are permanently blocked (walls) and therefore
unreachable.  The blocked flat indices are $\{9, 12, 17, 20, 25, 28, 29, 30,
33, 34, 38, 42, 46, 50, 54\}$.  The remaining 48 non-terminal, non-blocked
cells form the effective state space.

\paragraph{Action space.}
Nine deterministic intended-move actions: N, NE, E, SE, S, SW, W, NW, and
Stay (action indices 0--8).  The \emph{intended} move is always attempted
first; the actual outcome is stochastic due to drift (see below).

\paragraph{Transition kernel.}
After the intended move resolves, a stochastic drift is applied independently.
With probability 0.70 no drift occurs; with probability 0.05 each the agent
drifts one step in each cardinal direction; and with probability 0.025 each
it drifts diagonally.  If the intended move or drift would take the agent
out-of-bounds or into a blocked cell, the move is \emph{reverted} to the
pre-move position.

\paragraph{Cost structure.}
Each step incurs a base cost of $1 + \mathcal{N}(0, 0.04)$ (standard
deviation 0.2).  If the intended move is blocked by a wall, an additional
penalty of $1000 + \mathcal{N}(0, 0.04)$ is added, creating a steep incentive
to avoid walls.  Drift that results in an actual displacement from the
post-action position incurs an extra cost of $1 + \mathcal{N}(0, 0.04)$.

\paragraph{Terminal state.}
Flat index 36 (row~4, column~4) is the goal.  It is absorbing: once reached,
the episode ends with zero additional cost.

\paragraph{Discount factor.}
$\gamma = 0.85$.  This moderately discounts future costs, balancing short-term
efficiency against long-term path quality.

\subsection{Objective}

The control objective is to find a stationary policy
$\pi^*: \mathcal{S} \to \mathcal{A}$ that minimises the expected discounted
cumulative cost from every starting state:
\begin{equation}
    V^{\pi^*}(s) = \min_{\pi}\; \mathbb{E}_{\pi}\!\left[
        \sum_{t=0}^{\infty} \gamma^t\, c_t \;\Big|\; s_0 = s
    \right], \quad \forall\, s \in \mathcal{S}.
\end{equation}
The known optimal values $V^*$ (computed offline via value iteration on the
model matrices $P$ and $C$) serve as ground truth for all evaluation metrics.
Values range from $V^*(s)\approx 1.0$ for states adjacent to the goal up to
$V^*(s)\approx 5.4$ for far corners.

% ============================================================
\section{Algorithms Implemented}
% ============================================================

All nine algorithms share a common interface: \texttt{select\_action(s)},
\texttt{update(s,\,a,\,c,\,s')}, \texttt{get\_value\_estimate()} and
\texttt{get\_policy()}.  State is always the flat integer index.  The
count-based learning rate $\alpha_t = n_t(s,a)^{-\omega}$ (with $\omega=0.8$,
counts reset each episode and initialised to 1) is shared across all TD
algorithms; policy-gradient methods use a fixed learning rate instead.

\subsection{Q-Learning}

Q-learning is an off-policy TD method.  At each step it updates:
\begin{equation}
    Q(s,a) \;\leftarrow\; Q(s,a) + \alpha\bigl[c + \gamma\min_{a'}Q(s',a') - Q(s,a)\bigr].
\end{equation}
The greedy bootstrap $\min_{a'}Q(s',a')$ decouples the target from the
behaviour policy, allowing convergence to $Q^*$ under mild conditions.
Exploration uses $\varepsilon$-greedy with $\varepsilon_0=0.4$, decaying
multiplicatively by $0.9$ per step within each episode.

\textbf{Hyperparameter tuning.}  The decay factor 0.9 was chosen to ensure
$\varepsilon$ collapses quickly within each 300-step episode (reaching
$\approx 0.4 \times 0.9^{50} \approx 0.003$ after 50 steps), promoting
exploitation after the first few steps.  $\varepsilon_0=0.4$ balances initial
exploration without wasting too many early steps on random wall collisions.

\subsection{Double Q-Learning}

Double Q-learning maintains two tables $Q_A$ and $Q_B$.
At each step, one table is chosen with probability 0.5 to be updated; the
other evaluates the target:
\begin{equation}
    \text{target} = c + \gamma\, Q_B\!\bigl(s', \arg\min_{a'} Q_A(s',a')\bigr).
\end{equation}
This decouples \emph{action selection} from \emph{action evaluation},
reducing the maximisation bias present in standard Q-learning.  Exploration
uses the mean $(Q_A+Q_B)/2$.

\textbf{Hyperparameter tuning.}  Separate visit counts per table are maintained
so each table's learning rate adapts only to how often it was updated,
preventing one table from ageing prematurely.

\subsection{SARSA}

SARSA is an on-policy TD method.  The next action
$a'$ is sampled from the current $\varepsilon$-greedy policy before the update:
\begin{equation}
    Q(s,a) \;\leftarrow\; Q(s,a) + \alpha\bigl[c + \gamma\, Q(s',a') - Q(s,a)\bigr].
\end{equation}
Because the target uses the same policy that generates behaviour, SARSA
optimises for the \emph{behaviour} policy and is inherently cautious near
dangerous states (walls).

\textbf{Hyperparameter tuning.}  Same $\varepsilon$ schedule as Q-learning.
The on-policy nature means no additional tuning is needed; the conservative
tendency is an emergent structural property.

\subsection{Expected SARSA}

Expected SARSA replaces the sampled $Q(s', a')$ with its expectation under
the current $\varepsilon$-greedy policy:
\begin{equation}
    \text{target} = c + \gamma \sum_{a'} \pi(a'|s')\, Q(s',a').
\end{equation}
This eliminates one source of variance compared to SARSA (no random sample
of $a'$) while remaining on-policy.  If the limit $\varepsilon \to 0$ it
reduces to Q-learning.

\textbf{Hyperparameter tuning.}  Identical schedule to SARSA; the variance
reduction is structural.

\subsection{Entropy-Regularised Q-Learning (G-Learning)}

G-learning augments the Bellman equation with a KL
divergence penalty relative to a uniform prior $\rho(a|s)=1/9$, replacing
the hard-$\min$ bootstrap with a softmin value:
\begin{equation}
    V(s') = -\frac{1}{\beta}\log\!\left(\sum_{a'} e^{-\beta\, G(s',a')}\right)
    \quad(\text{with min-subtraction for stability}).
\end{equation}
The induced policy is $\mu(a|s) \propto \exp(-\beta\, G(s,a))$.  The inverse
temperature $\beta$ follows a linear schedule:
\begin{equation}
    \beta_t = \beta_{\min} + t\,\frac{\beta_{\max}-\beta_{\min}}{N},\quad
    \beta_{\min}=0.1,\quad \beta_{\max}=7.0.
\end{equation}

\textbf{Hyperparameter tuning.}  $\beta_{\min}=0.1$ produces a near-uniform
softmin policy early in training (maximum exploration); $\beta_{\max}=7.0$
approximates the greedy policy at convergence.  The linear schedule over 1500
episodes ensures a smooth transition.

\subsection{REINFORCE}

REINFORCE is a Monte-Carlo policy gradient method.  A full episode is collected, the discounted return $G_t = \sum_{k=t}^{T}\gamma^{k-t}c_k$
is computed at each step, and a softmax policy $\pi(a|s)=\text{softmax}(\theta_s)$
is updated by gradient descent on the expected cost:
\begin{equation}
    \theta_s \;\leftarrow\; \theta_s + \alpha\, G_t\,(\pi(\cdot|s) - \mathbf{e}_{a_t}).
\end{equation}

\textbf{Hyperparameter tuning.}  Fixed $\alpha=0.01$ for the policy.  A
smaller rate would slow convergence; a larger one risks instability due to the
high-variance returns produced by wall collisions.

\subsection{REINFORCE with Baseline}

A learned state-value baseline $V(s)$ is subtracted from $G_t$ to reduce
gradient variance:
\begin{align}
    \delta_t &= G_t - V(s_t),\\
    V(s_t) &\;\leftarrow\; V(s_t) + \alpha_b\,\delta_t,\\
    \theta_s &\;\leftarrow\; \theta_s + \alpha\,\delta_t\,(\pi(\cdot|s) - \mathbf{e}_{a_t}).
\end{align}
The baseline does not introduce bias because $V(s_t)$ is independent of the
action taken.

\textbf{Hyperparameter tuning.}  $\alpha_b = 0.1$ (baseline), $\alpha = 0.01$
(policy).  A larger $\alpha_b$ allows the baseline to track the mean return
quickly, reducing variance sooner.

\subsection{One-Step Actor-Critic}

The one-step actor-critic updates online at every step using the TD error
$\delta = c + \gamma V(s') - V(s)$ as the learning signal for both critic
and actor:
\begin{align}
    V(s) &\;\leftarrow\; V(s) + \alpha_c\,\delta,\\
    \theta_s &\;\leftarrow\; \theta_s + \alpha_a\,\delta\,(\pi(\cdot|s) - \mathbf{e}_{a_t}).
\end{align}

\textbf{Hyperparameter tuning.}  $\alpha_c = \alpha_a = 0.01$.  Online updates
reduce the per-episode variance compared to REINFORCE at the cost of some bias.

\subsection{Actor-Critic with Eligibility Traces (AC($\lambda$))}

AC($\lambda$) maintains eligibility trace vectors $\mathbf{z}_V$ (shape 64)
and $\mathbf{z}_\theta$ (shape $64\times9$).  At each step:
\begin{align}
    \mathbf{z}_V &\;\leftarrow\; \gamma\lambda_c\,\mathbf{z}_V + \mathbf{e}_s,\\
    \mathbf{z}_\theta &\;\leftarrow\; \gamma\lambda_a\,\mathbf{z}_\theta
        + (\pi(\cdot|s)-\mathbf{e}_{a_t})\,\text{at row }s,\\
    V &\;\leftarrow\; V + \alpha_c\,\delta\,\mathbf{z}_V,\qquad
    \theta \;\leftarrow\; \theta + \alpha_a\,\delta\,\mathbf{z}_\theta.
\end{align}
Traces are reset to zero at the start of each episode.

\textbf{Hyperparameter tuning.}  $\lambda_c = \lambda_a = 0.9$ interpolates
between one-step TD ($\lambda=0$) and Monte-Carlo ($\lambda=1$), providing a
good bias-variance trade-off in a 300-step episode.  The trace decay was
confirmed to be applied \emph{globally} across all states at each step,
a correction identified during debugging (see Section~\ref{sec:llm}).

% ============================================================
\section{Experimental Setup}
% ============================================================

All experiments are governed by a single configuration file
(\texttt{experiments/configs/default.yaml}) to ensure no algorithm benefits
from private hyperparameter tuning.

\begin{table}[H]
\centering
\caption{Shared experimental settings.}
\label{tab:setup}
\begin{tabular}{ll}
\toprule
Parameter & Value \\
\midrule
Discount factor $\gamma$ & 0.85 \\
Training episodes $N$ & 1500 \\
Max steps per episode & 300 \\
Independent runs (seeds) & 10 \\
LR exponent $\omega$ & 0.8 \\
$\varepsilon_0$ (TD methods) & 0.4 \\
$\varepsilon$ decay per step & 0.9 \\
\bottomrule
\end{tabular}
\end{table}

Seeds are assigned as $\text{seed}_\text{run} = 1000r$ ($r=0,\ldots,9$) for
the agent and $\text{seed}+1$ for the environment, ensuring the two are
independently seeded but reproducible.  All random operations go through a
seeded \texttt{numpy.random.Generator} (never the global state).

\subsection{Performance Metrics}

Three metrics are recorded after each episode using the agent's current
value estimate $\hat{V}(s)$ and induced policy $\hat{\pi}$:

\begin{enumerate}[leftmargin=*]
    \item \textbf{Signed relative error} (value bias):
        \[
            e_\pm = \frac{1}{|\mathcal{S}_v|}
            \sum_{s\in\mathcal{S}_v} \frac{\hat{V}(s) - V^*(s)}{V^*(s)},
        \]
        where $\mathcal{S}_v$ excludes blocked and terminal states.
        Positive values indicate over-estimation (optimism); negative values
        indicate under-estimation (pessimism).
    \item \textbf{Absolute relative error} (value accuracy):
        \[
            e_\text{abs} = \frac{1}{|\mathcal{S}_v|}
            \sum_{s\in\mathcal{S}_v} \frac{|\hat{V}(s) - V^*(s)|}{V^*(s)}.
        \]
    \item \textbf{Policy evaluation error} (primary metric, computed every 50
        episodes):
        \[
            e_\pi = \frac{1}{|\mathcal{S}_v|}
            \sum_{s\in\mathcal{S}_v} \frac{|V^{\hat{\pi}}(s) - V^*(s)|}{V^*(s)},
        \]
        where $V^{\hat{\pi}}$ is computed by model-based policy evaluation
        of the greedy (or softmax) policy induced by the current agent.
\end{enumerate}

The policy evaluation error is the primary cross-family comparison metric
because $e_\text{abs}$ has incompatible semantics across algorithm families:
for TD methods $\hat{V}(s)=\min_a Q(s,a)\approx V^*(s)$, while for
actor-critic and REINFORCE agents $\hat{V}(s)=V^{\hat{\pi}}(s)$, which
starts very high (a random policy hitting walls repeatedly produces
$V^{\hat{\pi}}\gg V^*$).

\subsection{Fairness of Comparison}

\begin{itemize}[leftmargin=*]
    \item \textbf{Same environment instances.}  Per-run, both the agent and the
        environment share the same base seed, so every algorithm faces
        identical stochastic conditions within a run.
    \item \textbf{Same number of interactions.}  All algorithms train for
        exactly 1500 episodes with at most 300 steps each; no algorithm
        receives more environment samples.
    \item \textbf{Same discount factor.}  $\gamma=0.85$ for all.
    \item \textbf{Adaptive LR for TD methods.}  The count-based schedule
        $\alpha_t = n_t^{-0.8}$ is applied uniformly to Q-learning, Double Q,
        SARSA, Expected SARSA, and G-learning, rather than hand-tuning per
        algorithm.
    \item \textbf{Fixed LR for PG/AC methods.}  REINFORCE and AC use fixed
        $\alpha=0.01$, consistent with gradient-descent practice where
        count-based rates are less natural.
    \item \textbf{No shared state.}  Each run creates a fresh agent and
        environment instance; no gradient or Q-table information is shared
        across seeds.
\end{itemize}

% ============================================================
\section{Results}
% ============================================================

\subsection{Learning Curves}

\begin{figure}[H]
    \centering
    \includegraphics[width=\linewidth]{../report/figures/05_policy_err_all.png}
    \caption{Policy evaluation error over training episodes for all nine
    algorithms.  Lower is better.  Shaded bands show $\pm 1$ standard
    deviation across 10 runs.  Plotted every 50 episodes (30 checkpoints
    total).}
    \label{fig:policy_err_all}
\end{figure}

\begin{figure}[H]
    \centering
    \begin{subfigure}[b]{0.48\linewidth}
        \includegraphics[width=\linewidth]{../report/figures/05_abs_err_grouped.png}
        \caption{Absolute relative error, grouped by algorithm family.}
        \label{fig:abs_err_grouped}
    \end{subfigure}
    \hfill
    \begin{subfigure}[b]{0.48\linewidth}
        \includegraphics[width=\linewidth]{../report/figures/05_signed_err_td.png}
        \caption{Signed relative error for TD algorithms, showing the
        optimism/pessimism dynamics.}
        \label{fig:signed_err}
    \end{subfigure}
    \caption{Value-estimate quality metrics for TD-based algorithms.}
\end{figure}

\begin{figure}[H]
    \centering
    \includegraphics[width=0.85\linewidth]{../report/figures/05_summary_bar.png}
    \caption{Final (episode 1500) mean policy evaluation error across 10 runs.
    Error bars show one standard deviation.}
    \label{fig:summary_bar}
\end{figure}

\begin{figure}[H]
    \centering
    \includegraphics[width=\linewidth]{../report/figures/05_value_heatmaps.png}
    \caption{Learned value-function heatmaps (single seed, 1500 episodes) for
    all nine algorithms alongside the optimal $V^*$.  Warmer colours indicate
    higher cost-to-go.  Blocked and terminal cells are masked.}
    \label{fig:heatmaps}
\end{figure}

\begin{figure}[H]
    \centering
    \includegraphics[width=\linewidth]{../report/figures/05_policy_grids.png}
    \caption{Learned policies visualised as action arrows.  Each arrow points
    in the direction of the highest-probability action at each valid state.}
    \label{fig:policies}
\end{figure}

Table~\ref{tab:final_results} summarises the final-episode policy evaluation
errors from a single representative seed (seed 42, 1500 episodes), which is
consistent with multi-run trends.

\begin{table}[H]
\centering
\caption{Final policy evaluation error after 1500 training episodes (seed 42).
Lower is better.}
\label{tab:final_results}
\begin{tabular}{lcc}
\toprule
Algorithm & Family & Policy Eval Error \\
\midrule
Q-learning            & Off-policy TD      & \textbf{0.0447} \\
G-learning            & Off-policy TD (ER) & 0.1259 \\
Double Q-learning     & Off-policy TD      & 0.1779 \\
Expected SARSA        & On-policy TD       & 0.3895 \\
AC($\lambda$)         & Actor-Critic       & 0.5130 \\
One-step Actor-Critic & Actor-Critic       & 0.5830 \\
REINFORCE+baseline    & Policy Gradient    & 0.6038 \\
SARSA                 & On-policy TD       & 0.7490 \\
REINFORCE             & Policy Gradient    & 0.9550 \\
\bottomrule
\end{tabular}
\end{table}

\subsection{Final Conclusion and Discussion}

\paragraph{Q-learning is the best performer for our case.}
Q-learning achieves the lowest final policy evaluation error ($e_\pi \approx
0.045$).  The key reasons are structural.  First, the off-policy TD target
$c + \gamma\min_{a'}Q(s',a')$ bootstraps directly toward $Q^*$ regardless
of the exploration policy; the agent does not need to follow the optimal
policy to learn it.  Second, the count-based learning rate $\alpha_t =
n_t^{-0.8}$ satisfies the Robbins--Monro conditions ($\sum\alpha_t=\infty$,
$\sum\alpha_t^2<\infty$) and adapts per state-action pair, giving more weight
to early visits (rapid initial learning) while slowing down once a cell is
well-characterised.  Third, Q-learning updates every step, accumulating
1500 $\times$ (up to 300) individual Bellman updates, whereas REINFORCE
updates only once per episode.

\paragraph{G-learning is the best-performing novel algorithm.}
G-learning (entropy-regularised Q) achieves $e_\pi \approx 0.126$, second
only to Q-learning.  The entropy bonus encourages the agent to keep its
policy spread across actions early in training (when $\beta$ is small), which
naturally provides exploration without relying solely on $\varepsilon$-greedy.
As $\beta$ increases linearly over 1500 episodes, the softmin policy sharpens
toward the greedy policy.  This annealing is principled: it corresponds to
minimising a KL-regularised cost, so the intermediate ``soft'' policies are
optimal for their respective temperature.  In environments with high wall
penalties, maintaining some probability on all actions early prevents the
agent from prematurely committing to paths that lead to walls.

\paragraph{Double Q-learning partially mitigates maximisation bias.}
Standard Q-learning is known to overestimate action values because the same
sample is used to both select and evaluate the greedy action.  In a
cost-minimisation setting this manifests as \emph{pessimism}: $\min Q$ tends
to underestimate true costs early in training, which can be read from the
signed error plots (negative $e_\pm$ early on for Q-learning, closer to zero
for Double Q).  Double Q corrects this, and its final error of 0.178 is
moderately better than Expected SARSA (0.390) but slightly worse than Q-learning,
partly because each table receives only half as many updates per episode.

\paragraph{On-policy TD methods learn more cautiously.}
SARSA's final policy error (0.749) is substantially larger than Q-learning's
(0.045) on the same seed.  The theoretical reason is that SARSA optimises
$Q^\pi$ for the \emph{behaviour} policy $\pi$ (epsilon-greedy) rather than
the optimal policy.  In a wall-heavy environment, the epsilon-greedy policy
occasionally takes random actions near walls, incurring the 1000-unit penalty.
The SARSA target therefore includes these costly exploratory outcomes,
making the learned $Q^\pi$ overestimate the cost of near-wall states and
leading to an overly conservative policy that avoids regions near the goal
boundaries even after epsilon has decayed.

Expected SARSA (0.390) significantly outperforms SARSA by replacing the
stochastic sample $Q(s',a')$ with the analytical expectation under $\pi$.
This variance reduction accelerates convergence: fewer noisy target evaluations
mean cleaner Q-table updates, especially early in training when $\varepsilon$
is large.

\paragraph{Effect of exploration.}
The epsilon-greedy scheme decays per step within each episode (not per episode
globally).  With $\varepsilon_0=0.4$ and decay 0.9, epsilon reaches $\approx 0.003$
after 50 steps.  This means each episode transitions from exploration to
exploitation within its first 50 steps, which is appropriate for episodes
capped at 300 steps.  Insufficient exploration would leave large parts of the
grid unvisited; excessive exploration would accumulate too many wall penalties.
G-learning's entropy bonus provides a richer form of directed stochasticity than
hard epsilon-greedy: the softmin policy samples actions proportionally to their
estimated optimality, rather than either purely greedy or uniformly random.

\paragraph{Entropy regularisation helps here.}
G-learning outperforms standard Q-learning variants from other families and
SARSA variants.  The KL penalty toward a uniform prior effectively prevents
premature certainty, especially valuable in the early episodes when $Q$-values
are far from their true values.  Because the Gridworld has a clear bottleneck
structure (states near blocked columns must choose specific directions), the
uniform prior's pull ensures the agent explores those corridors, rather than
fixating on a locally cheap path discovered by chance.

\paragraph{Policy-based methods are the worst performers.}
Pure REINFORCE achieves the worst policy error (0.955).  The fundamental issue
is \emph{variance}: the return $G_t = \sum_{k\geq t}\gamma^k c_k$ is a sum
over the entire remaining episode.  A single wall collision ($c=1001$) can
dominate $G_t$ for all preceding steps, producing an enormous and misleading
gradient signal.  With 300 steps per episode and Gaussian noise on each cost,
the variance of $G_t$ is enormous.  REINFORCE also does not update its policy
between steps, losing the per-step Bellman information that TD methods exploit.

Adding a baseline (REINFORCE+baseline) reduces the error to 0.604 by centering
the gradient signal: $\delta_t = G_t - V(s_t)$ is smaller in magnitude than
$G_t$ alone once $V$ has learned a reasonable estimate of $V^\pi$.  The
variance reduction is especially pronounced near the goal, where $V(s)$ is
accurate and $\delta_t \approx 0$ for optimal steps.

\paragraph{Actor-Critic methods offer a middle ground.}
One-step AC ($e_\pi=0.583$) outperforms pure REINFORCE ($e_\pi=0.955$) because
it updates online: the per-step TD error $\delta = c + \gamma V(s') - V(s)$ is
a much lower-variance signal than a full Monte-Carlo return.  However, it is
biased: bootstrapping from a value function that is itself being learned
introduces approximation error that only vanishes at convergence.

AC($\lambda$) with $\lambda=0.9$ achieves $e_\pi=0.513$, better than one-step
AC.  Eligibility traces interpolate between one-step TD and MC returns: at
$\lambda=0.9$, the effective return is a $\lambda$-weighted combination, giving
the algorithm a longer temporal horizon.  This helps credit assignment: in
navigating a maze, the reward (reaching the terminal state) is temporally
distant from the actions that made it possible; one-step TD assigns credit only
to the last action before arrival, while traces propagate credit backward to
earlier steps that set up the winning trajectory.

\paragraph{Hyperparameter sensitivity.}
TD methods with count-based learning rates are robust: the rate adapts
automatically and requires only the exponent $\omega=0.8$ to be specified.
Policy-gradient and actor-critic methods are more sensitive to the fixed
learning rate $\alpha=0.01$.  Too large a value causes oscillation when
early episodes include extreme wall penalties; too small a value slows
convergence.  The choice 0.01 is conservative and proven stable across all
seeds, but a larger value (e.g.\ 0.05) might accelerate early convergence
at the risk of instability.

\paragraph{Summary.}
Off-policy value-based methods—especially Q-learning—are most effective on
this tabular, cost-minimisation Gridworld because they exploit the Bellman
structure directly at every step without on-policy bias or Monte-Carlo
variance.  On-policy TD methods converge more slowly due to conservative
target evaluation.  Policy-gradient methods struggle because the high-variance
cost signal (wall penalties up to 1001) overwhelms gradient estimates
before the policy has had time to learn to avoid walls.  Entropy regularisation
(G-learning) is the most effective modification within the TD family, providing
principled exploration annealing.  Eligibility traces in AC($\lambda$) provide
a meaningful improvement over one-step AC by addressing the temporal
credit-assignment problem.

% ============================================================
\section{Use of LLMs}
\label{sec:llm}
% ============================================================

LLMs (Claude, via Anthropic's Claude Code CLI) were used throughout the project
for implementation and debugging.  The LLM usage log is maintained in
\texttt{report/llm\_usage.md}; key usage sessions are summarised below.

\begin{enumerate}[leftmargin=*]
    \item \textbf{Environment porting (2026-04-02):} The \texttt{GridWorldEnv}
        class, transition matrices $P$ and $C$, and \texttt{policy\_eval} were
        ported from a reference notebook into a structured module following the
        project interface.  The only intentional deviation was replacing global
        \texttt{np.random} calls with a seeded \texttt{np.random.Generator} for
        reproducibility.  Step dynamics and matrix logic were verified to match
        the notebook exactly.

    \item \textbf{Project scaffolding (2026-04-02):} The full directory structure,
        abstract base class, stub files, configuration, and test skeletons were
        generated.  All decisions (e.g.\ count initialisation to 1, per-episode
        reset) were derived from the specification and verified against the
        reference notebook.

    \item \textbf{Algorithm implementation (2026-04-02 and 2026-04-06):}
        All nine algorithm files and their corresponding test classes
        (48 tests total) were implemented.  The implementations followed
        the mathematical specifications in the project brief; the LLM helped
        translate equations into vectorised NumPy code.  Each implementation
        was verified on a trivial 2-state 2-action MDP with known $Q^*$ before
        being run on the full Gridworld.

    \item \textbf{AC($\lambda$) bugfix and experiment runners (2026-04-23):}
        Two bugs in AC($\lambda$) were identified and fixed: (a) eligibility
        traces were only decayed at the current state, not globally, causing
        unbounded growth; (b) the actor trace used the reward-maximisation sign
        convention instead of cost-minimisation.  The LLM generated a candidate
        fix that was verified by observing policy error decreasing (from stuck
        at 1.50 to a trajectory toward 0.51) after the correction.

    \item \textbf{Report writing:} LLM used to assist in drafting, structuring,
        and refining the technical content of this report.  All theoretical
        interpretations and experimental conclusions were verified by the group
        against the observed numerical results and course material.
\end{enumerate}

The final implementation, experiment design, numerical results, and conclusions
were verified and finalised by the group.

% ============================================================
\section{Division of Work}
% ============================================================

\begin{itemize}[leftmargin=*]
    \item \textbf{Nirvan Nagrale (2022ME12048):} Environment implementation and
        verification; Q-learning, Double Q-learning, and G-learning
        implementations; experiment runner (\texttt{run\_all.py}); debugging of
        AC($\lambda$); report writing (Sections 2, 3, 5).

    \item \textbf{Kunal Parghi (2022ME12043):} SARSA and Expected SARSA
        implementations; unit test suite; hyperparameter tuning experiments;
        report writing (Sections 1, 4, 6).

    \item \textbf{Aditya Yadav (2022ME11320):} REINFORCE, REINFORCE+baseline,
        one-step Actor-Critic, and AC($\lambda$) implementations; plotting
        utilities; notebook analysis and figure generation; report writing
        (Section 5 discussion, abstract).
\end{itemize}

All members participated in code review, experimental validation, and the
final report revision.

% ============================================================
\section*{References}
% ============================================================

\begin{enumerate}[leftmargin=*]
    \item R.\ Fox, A.\ Pakman, N.\ Tishby. \textit{Taming the Noise in
        Reinforcement Learning via Soft Updates.} UAI 2016.

    \item C.\ J.\ C.\ H.\ Watkins, P.\ Dayan. \textit{Q-learning.}
        Machine Learning, 8(3-4):279--292, 1992.

    \item H.\ van Hasselt. \textit{Double Q-learning.}
        Advances in Neural Information Processing Systems, 2010.

    \item G.\ A.\ Rummery, M.\ Niranjan. \textit{On-line Q-learning using
        connectionist systems.} Technical Report, Cambridge, 1994.

    \item R.\ J.\ Williams. \textit{Simple statistical gradient-following
        algorithms for connectionist reinforcement learning.}
        Machine Learning, 8(3-4):229--256, 1992.

    \item R.\ S.\ Sutton, A.\ G.\ Barto. \textit{Reinforcement Learning:
        An Introduction.} 2nd ed., MIT Press, 2018.
\end{enumerate}

\end{document}
